\subsection{La capa de trayectoria: del seguidor de puntos al juego continuo}

Las primeras corridas del banco movían cada coalición con un seguidor de
puntos de espera y frenados heurísticos; las trayectorias resultantes eran
erráticas —bucles frente al vano, zigzags de tracking— y se parecían más a un
reactivo que al juego de trayectorias de \cref{sec:continuacion}. Se
implementó entonces esa capa tal como está formulada allí, y el banco compara
ambas.

Cada jugador —una coalición en transporte o un AMR libre— mantiene una
continuación de $N$ nodos separados $\Delta$ ($N=12$, $\Delta=0{,}5$~s para
coaliciones; $N=10$, $\Delta=0{,}4$~s para AMR) y cada $0{,}2$~s la deforma
por gradiente proyectado del coste \eqref{eq:coste-trayectoria} en su forma
implementada: atracción terminal y de progreso hacia el siguiente punto de la
ruta, suavidad (aceleración discreta), esfuerzo (longitud recorrida, para que
ceder signifique frenar y no dar vueltas), límite físico de velocidad, barrera
frente a paredes con la holgura lateral real de la carga alineada, barrera
frente a cargas ajenas, y la externalidad atómica exacta con las
continuaciones publicadas por los demás jugadores en los mismos instantes
(\cref{th:externalidad-atomica}), re-muestreadas cuando los horizontes
difieren. La única propuesta inicial es la recta al objetivo; el arranque en
caliente desplaza la continuación conforme avanza la planta; el primer
segmento de la continuación es la velocidad comandada al mercado de
\emph{wrench}. Cada publicación de continuación cuenta como mensaje a los
vecinos de $n$ saltos. La ruta discreta que la homotopía impone
(\cref{obs:homotopia}) es explícita: carril de espera lateral con circulación
por la derecha, eje del vano, carril de salida; quien espera el arrendamiento
publica su posición como continuación estática, para que la propietaria no le
ceda el paso a un plan que no va a ejecutarse —la cortesía mutua produjo
interbloqueo en la primera versión—. Los AMR ocupados en formación o acoplados
publican también su posición, y los libres planifican contra ellos, contra las
cargas y contra las paredes.

Dos límites de esta implementación se declaran. El gradiente es local: un
objetivo situado tras un muro hace que la continuación resbale a lo largo de
la pared, de ahí la ruta por puntos; y las interacciones por pares con radios
envolventes producen desvíos amplios cuando el carril de espera de una
coalición cae sobre el camino de salida de la otra, que es lo que se ve en
\cref{fig:sim-plan}. Lo que el trazo ejecutado conserva de ondulación no viene
del plan sino del lazo físico —tres AMR que giran bajo la plataforma para
producir fuerza lateral— y se redujo con amortiguamiento de rumbo y de guiñada.
